\documentclass[11pt]{article}

\usepackage[letterpaper,left=1in,right=1in,top=0.75in,bottom=0.75in]{geometry}
\usepackage{microtype}
\usepackage{amsmath}
\usepackage{amssymb}
\usepackage{booktabs}
\usepackage{array}
\usepackage{tabularx}
\usepackage[table,dvipsnames]{xcolor}
\usepackage{enumitem}
\usepackage[most]{tcolorbox}
\usepackage{float}
\usepackage[hidelinks]{hyperref}

\emergencystretch=3em

\definecolor{infoblue}{HTML}{2C3E50}
\definecolor{rule}{HTML}{BDC3C7}
\definecolor{decisiongreen}{HTML}{1B7F5A}
\definecolor{rejectred}{HTML}{A33A2B}

\newcommand{\code}[1]{\texttt{#1}}
\newcommand{\rfckw}[1]{\textsc{#1}}
\newcommand{\origpaper}{\href{https://stackarmor.github.io/rfc-fedramp-vdr/vdr-pain-cvss.pdf}{the initial PAIN publication}}
\newcommand{\ceilingpaper}{\href{https://stackarmor.github.io/rfc-fedramp-vdr/vdr-pain-security-requirements.pdf}{\emph{Before the PAIN Equation}}}
\newcolumntype{P}[1]{>{\raggedright\arraybackslash}p{#1}}
\newcolumntype{Y}{>{\raggedright\arraybackslash}X}

\newtcolorbox{plain}{
  enhanced, breakable,
  colback=Sepia!7, colframe=Sepia!55, boxrule=0.4pt, arc=2pt,
  left=7pt, right=7pt, top=4pt, bottom=4pt,
  fontupper=\small,
  coltitle=Sepia!72!black, fonttitle=\bfseries\sffamily\footnotesize,
  title={IN PLAIN TERMS}
}

\newtcolorbox{decisionbox}{
  enhanced, breakable,
  colback=decisiongreen!5, colframe=decisiongreen!65, boxrule=0.5pt, arc=2pt,
  left=7pt, right=7pt, top=5pt, bottom=5pt,
  fontupper=\small,
  coltitle=decisiongreen!75!black, fonttitle=\bfseries\sffamily\footnotesize,
  title={FINAL CALIBRATION}
}

\newtcolorbox{cautionbox}{
  enhanced, breakable,
  colback=rejectred!4, colframe=rejectred!55, boxrule=0.4pt, arc=2pt,
  left=7pt, right=7pt, top=4pt, bottom=4pt,
  fontupper=\small,
  coltitle=rejectred!75!black, fonttitle=\bfseries\sffamily\footnotesize,
  title={IMPORTANT BOUNDARY}
}

\title{\bfseries Calibrating PAIN Without Abandoning CVSS:\\
High-Centered Normalization and\\
Standards-Anchored Thresholds}
\author{Matthew Venne \\ \small Chief Technology Officer, stackArmor}
\date{July 2026\\\small Version 1.1}

\begin{document}
\maketitle

\begin{abstract}
\noindent
The initial paper \emph{A Deterministic, CVSS--Environmental Method for FedRAMP
VDR/VER Vulnerability Prioritization} introduced a reproducible method for deriving
Potential Agency Impact (PAIN) from a vulnerability's confidentiality, integrity,
and availability impacts and an asset's CVSS Environmental Security Requirements.
Operational backtesting exposed a calibration defect: the method normalized against
the CVSS v3.1 Medium-requirement impact ceiling while High requirements could produce
a larger pre-cap aggregate. A single High technical impact on a High-requirement
dimension therefore normalized to approximately $0.918$ and crossed directly into
Debilitating. In one scan, approximately 40 percent of affected-resource observations
were N4.

This paper documents the alternatives evaluated and the resulting revision. More
asset sub-profiles could reduce severity but would multiply tagging choices. Lowering
the High requirement multiplier to $1.25$ approached the desired distribution but
discarded a CVSS-defined constant. Boolean component gates expressed the intended
compound-impact semantics but added a second classifier and discarded some continuous
dimensional information. Raising a threshold alone treated the symptom while leaving
the normalization axis centered on Medium requirements.

The selected method retains the CVSS impact weights, Security Requirement weights,
and complement-product aggregation. It replaces the saturated $0.915$ normalization
with the maximum attainable pre-cap aggregate under High requirements,
$D_H=0.995904$. Rather than carrying forward the original lower thresholds, Narrow
begins at the first complete technical loss on a Low-requirement dimension and
Disruptive begins at a High impact on a Moderate-requirement dimension. The
Debilitating threshold is set to $0.933$, which lies
strictly between the highest possible non-compound combination ($0.931993$) and the
lowest possible compound catastrophic combination ($0.933423$). Exhaustive testing
of all 729 discrete impact/requirement combinations and backtesting against two
anonymized evaluation datasets support the result. All three boundaries are therefore
anchored to stated consequence scenarios and empirically tested, not fitted to a
desired histogram.
\end{abstract}

\tableofcontents

\section{Status and relationship to the initial method}

This document is an informational calibration companion to \origpaper{}. It does
not define a FedRAMP requirement, change a CVSS score, or claim that the resulting
scalar is a conformant CVSS Environmental Score. It revises the PAIN-specific step
that transforms CVSS Environmental impact inputs into a unit scalar and then into
the words Minimal, Narrow, Disruptive, and Debilitating.

The initial method made three separations that remain intact:

\begin{itemize}[leftmargin=1.5em]
\item CVSS-derived impact and asset requirements determine the PAIN word.
\item Agency scope maps the word to N1--N5.
\item Exploitability and internet reachability select the remediation column; they
  do not change potential impact.
\end{itemize}

This calibration changes only the first of those steps. Remediation matrices,
likely-exploitable logic, internet-reachability logic, multi-agency handling, VEX
dispositions, and evidence-gated Modified impact reductions remain outside its
scope.

The governed asset-archetype catalog and any intended-use Security Requirements
ceiling are also external inputs to this calibration. This paper does not derive
either input. It holds the effective CR/IR/AR vector constant while comparing
scoring methods. \ceilingpaper{} specifies the upstream derivation that caps a
system-relative archetype vector by the confirmed federal information types in the
affected agency scope.

\begin{plain}
The original model asked the right question but used the wrong ruler at the top of
the scale. It measured High-requirement assets against a maximum created by Medium
requirements. This paper replaces the ruler, preserves the CVSS inputs, and places
the Debilitating line at the first combination that is both catastrophic and
compound.
\end{plain}

\section{The calibration problem}
\label{sec:problem}

Let $C,I,A$ be the effective CVSS impact values for confidentiality, integrity,
and availability:
\[
C,I,A \in \{\,0,\ 0.22,\ 0.56\,\},
\]
and let the corresponding asset requirements be
\[
\mathrm{CR},\mathrm{IR},\mathrm{AR}
\in \{\,0.5,\ 1.0,\ 1.5\,\}.
\]
These are the CVSS v3.1 None/Low/High impact weights and Low/Medium/High Security
Requirement weights. CVSS v3.1 defines the capped Modified Impact Sub-Score as
\begin{equation}
\mathrm{MISS}=\min\!\left(
1-(1-C\mathrm{CR})(1-I\mathrm{IR})(1-A\mathrm{AR}),\ 0.915
\right).
\label{eq:miss}
\end{equation}

The initial PAIN method normalized this value:
\begin{equation}
S_M=\frac{\mathrm{MISS}}{0.915},
\label{eq:oldscore}
\end{equation}
then used $(0.25,0.55,0.80)$ as the three word thresholds.

The value $0.915$ is the rounded Medium-requirement all-High impact ceiling:
\[
1-(1-0.56)^3=0.914816\approx0.915.
\]
But one High impact on a High requirement contributes $0.56\times1.5=0.84$.
Therefore a one-dimensional finding produces
\[
S_M=\frac{0.84}{0.915}=0.918033,
\]
which exceeds the original Debilitating threshold of $0.80$. This is the
single-dimension cliff: a complete loss of one important security property becomes
N4 on a single-agency asset before any second security property is affected.

The effect is operational, not cosmetic. For a Class C provider, an N4-to-N3
change moves the \code{LEV+NIRV} deadline from 8 days to 32 days and the
\code{LEV+IRV} deadline from 4 days to 16 days. A large cluster at one arithmetic
point can therefore create both remediation pressure and a foreseeable incentive
to assign a genuinely High requirement as Moderate.

\section{Design criteria}

The review applied five criteria before accepting a change:

\begin{enumerate}[leftmargin=1.7em]
\item \textbf{Semantic first.} Each PAIN boundary must have a written meaning that
  can be tested with scenarios. A target distribution is validation evidence, not
  the definition.
\item \textbf{Preserve CVSS-derived inputs.} The official $0.56/0.22/0$ impact
  weights, $1.5/1.0/0.5$ Security Requirement weights, and complement-product form
  should remain unless evidence demonstrates that a particular input is unsuitable.
\item \textbf{Keep tagging usable.} A regular operator should select an asset's role
  from a short governed catalog without navigating a cross-product of profiles,
  lenses, deployment modes, and impact exceptions.
\item \textbf{Preserve dimensional information.} Low impacts, Moderate requirements,
  and accumulation across C/I/A should continue to affect the scalar rather than
  disappear behind a boolean count.
\item \textbf{Remain deterministic and auditable.} Given the same effective impact
  vector and archetype, every implementation must produce the same result, retain
  full-precision intermediate values, and expose the policy version.
\end{enumerate}

\section{Alternatives evaluated}
\label{sec:alternatives}

\subsection{Additional asset sub-profiles}

One approach was to retain broad archetypes but add environment-specific variants,
such as internal-only orchestration, administrator-only CI/CD, redundant control
planes, or restricted security tooling. A primary archetype label could have been
combined with a second profile label, or the catalog could have introduced parallel
archetypes for each context.

Reviewing the underlying archetype taxonomy was still necessary: some broad High
assignments warranted clearer, role-based treatment. But taxonomy correction and
scoring calibration solve different problems. Reclassification is justified only
when the asset's actual mission role supports a different CR/IR/AR vector. Even a
well-governed catalog will contain genuine High requirements, and a scoring method
that turns every single H/H alignment into Debilitating remains miscalibrated.
Relying exclusively on re-categorizing assets would therefore move the calibration
problem into operator tagging, reward under-classification, and leave the underlying
arithmetic unchanged.

\paragraph{Illustrative profile model.}
The two-label form would be straightforward at first:

\begin{verbatim}
asset-archetype: orchestrator
asset-impact-profile: redundant-admin-only
\end{verbatim}

The second label would alter the requirement vector supplied by the first. Table~\ref{tab:profile-complexity}
shows representative outcomes considered during the design discussion. They are
illustrative rejected alternatives, not normative archetype assignments.

\begin{table}[H]
\centering
\footnotesize
\renewcommand{\arraystretch}{1.14}
\begin{tabularx}{\textwidth}{@{}P{0.19\textwidth}P{0.29\textwidth}cccY@{}}
\toprule
\textbf{Archetype} & \textbf{Impact profile} & \textbf{CR} & \textbf{IR} & \textbf{AR} & \textbf{Distinction encoded} \\
\midrule
\code{orchestrator} & service-critical control plane & M & M & H & Sustained loss stops service delivery. \\
\code{orchestrator} & redundant, administrator-only & M & M & M & Failover limits the availability consequence. \\
\code{cicd-pipeline} & production release/signing & H & H & M & Disclosure or corruption reaches production. \\
\code{cicd-pipeline} & isolated internal test & M & M & L & Output is non-production and reproducible. \\
\code{security-tooling} & inline preventive enforcement & M & H & H & Corruption or outage removes an active control. \\
\code{security-tooling} & passive telemetry only & M & M & L & Loss reduces visibility but not service delivery. \\
\code{identity-secrets} & durable credentials & H & M & M & Disclosure is catastrophic; loss is recoverable. \\
\code{identity-secrets} & replicated ephemeral token cache & H & M & L & Replication and regeneration reduce availability need. \\
\bottomrule
\end{tabularx}
\caption{Illustrative, non-normative sub-profile mappings considered and rejected.}
\label{tab:profile-complexity}
\end{table}

Even this small example requires rules for whether every profile is valid for every
archetype, what an omitted profile means, whether a profile replaces or modifies the
base vector, and whether a direct CR/IR/AR override takes precedence. Eight
archetypes and three profiles expose up to 24 apparent choices before exceptions;
adding a fourth profile adds another eight combinations and their documentation,
validation, tests, and migration behavior.

This could produce the desired numerical result because it moves more dimensions
from High to Moderate or Low before scoring. It was not selected for three reasons:

\begin{itemize}[leftmargin=1.5em]
\item Two independent metadata axes create a cross-product of valid and invalid
  combinations and make precedence rules necessary.
\item Parallel archetypes cause catalog growth, ambiguous choices, and drift as
  teams interpret nearly identical labels differently.
\item The approach moves calibration pressure into asset tagging. Operators can
  obtain a more favorable deadline by choosing a less severe profile even when the
  underlying system role is unchanged.
\end{itemize}

Environment-specific facts still matter, but they belong in evidence-backed CVSS
Modified metrics or other governed control evidence, not in an unlimited taxonomy
of scoring profiles.

\subsection{Changing the High multiplier from 1.50 to 1.25}

Reducing the High Security Requirement multiplier to $1.25$ directly removes the
single-dimension cliff:
\[
\frac{0.56\times1.25}{0.915}=0.765027,
\]
which is Disruptive at the original thresholds. Two High impacts on two High
requirements still produce
\[
\frac{1-(1-0.70)^2}{0.915}=0.994536,
\]
which remains Debilitating.

The distribution moved close to the desired range, but the alternative was not
selected. The $1.5$ value is the published CVSS v3.1 High requirement weight;
$1.25$ would be a PAIN-specific replacement without an independent derivation.
FIRST describes v3 calibration broadly as expert judgment checked against real
vulnerabilities, not as a public parameter-specific derivation of $1.5$. Retaining
it preserves a standardized input rather than claiming unique correctness for PAIN.
It would also alter every boundary, not only Debilitating. Finally, multiplier
reduction alone does not express the desired semantics: three High impacts on an
all-Moderate asset still produce $S\approx1$ and therefore N4.

\subsection{Raising the original Debilitating threshold}

The review was not opposed in principle to changing the $\theta$ values. Unlike the
CVSS metric weights, they are governed PAIN policy boundaries and should change when
the intended meaning of a PAIN word changes. The review did, however, reject
choosing thresholds primarily to manufacture a preferred distribution---for example,
tuning until Minimal, Narrow, Disruptive, and Debilitating appeared near the 20th,
40th, 60th, and 80th percentiles of the observed vulnerability population.

That approach would produce a tidy histogram but make classification relative to
the current dataset. The same impact and asset requirements could receive a
different PAIN word when unrelated vulnerabilities entered or left the population,
when a different tenant mix was scanned, or when asset-classification coverage
improved. Percentile targets may be useful diagnostics, but they are not defensible
definitions of consequence.

Moving $\theta_3$ from $0.80$ to $0.92$ while retaining the $0.915$ denominator
demotes the $0.918033$ single-dimension cluster. This is a legitimate PAIN policy
change, and the scan distribution improved. Standing alone, however, it retains a
Medium-centered axis and places the new line almost exactly above the cluster it is
intended to demote. Without a scenario-derived boundary, the change can reasonably
look like histogram fitting. The final method therefore changes $\theta_3$ only
after deriving a semantic boundary from the finite impact/requirement state space.
The lower thresholds are likewise re-derived from stated Minimal, Narrow, and
Disruptive anchor scenarios rather than inherited or percentile-fitted. The observed
histogram is used to test the consequence, not select any threshold.

\subsection{Strict two-dimension alignment gate}

A categorical alternative required two aligned pairs of High technical impact and
High asset requirement before N4:
\[
\#\{d:\ \mathrm{impact}_d=H\ \land\ \mathrm{requirement}_d=H\}\ge2.
\]
This made N4 rare, but it was too rigid. An archetype such as
\code{orchestrator=MMH} could never reach N4, even under complete
$C{:}H/I{:}H/A{:}H$ compromise, because only availability is High.

\subsection{Compound component gate}

The gate was broadened to require at least one catastrophic H/H alignment and at
least two High technical impacts on Moderate-or-High requirement dimensions. It
handled the orchestrator case and produced a credible distribution.

The method was not selected as the final implementation because it adds a second
classifier after the scalar and loses dimensional information at the boundary.
Low impacts cease to matter to eligibility, and vectors with materially different
continuous scores collapse to the same pair of counts. Implementations must also
explain why a raw scalar above the Debilitating threshold was subsequently capped
at Disruptive.

The gate nevertheless supplied the semantic test used to derive the final
arithmetic boundary.

\subsection{CVSS v4 MacroVectors}

CVSS v4 confirms that aligned impact and requirements can be treated categorically.
Its EQ6 group enters the more-severe level when \emph{any} of
\code{VC:H/CR:H}, \code{VI:H/IR:H}, or \code{VA:H/AR:H} is true, and it evaluates
EQ6 jointly with the vulnerable-system impact pattern in EQ3. CVSS v4 does not
require two aligned dimensions or compound impact. Its full score also includes
exploitability, subsequent-system impact, and threat metrics that PAIN deliberately
routes elsewhere.

Accordingly, CVSS v4 informed the alignment analysis but was not copied as the
PAIN classifier. The final method remains compatible with v4 vulnerable-system
impact values by mapping \code{VC/VI/VA} onto the PAIN C/I/A inputs when a v3.1
vector is unavailable.

\section{High-centered normalization}
\label{sec:normalization}

The final method separates the CVSS-defined cap from the underlying environmental
impact aggregate. Define the pre-cap aggregate
\begin{equation}
U=1-(1-C\mathrm{CR})(1-I\mathrm{IR})(1-A\mathrm{AR}).
\label{eq:aggregate}
\end{equation}

The symbol $U$ is PAIN notation for the expression before the CVSS cap; it is not a
new CVSS metric. CVSS v3.1 defines the conformant Modified Impact Sub-Score as
\begin{equation}
\mathrm{MISS}=\min(U,0.915).
\label{eq:miss-boundary}
\end{equation}
When $U\le0.915$, $U$ and MISS are numerically identical. When $U>0.915$, MISS
remains fixed at $0.915$ regardless of how much additional weighted impact the
other dimensions contribute. The cap belongs to the official CVSS 0--10 scoring
machinery, where MISS subsequently interacts with scope, exploitability, and
rounding rules.

PAIN does not reuse that downstream formula. Its severity axis is impact only;
exploitability and internet reachability are routed separately to remediation
timing. Consequently, collapsing every $U>0.915$ to the same value removes
information PAIN needs to distinguish different compound-impact states. The
revision branches immediately before the CVSS cap: the published CVSS calculation
continues to use MISS, while the PAIN calculation retains $U$.

The maximum attainable $U$ over the retained PAIN inputs occurs when all three
impacts and all three requirements are High:
\begin{equation}
D_H=1-(1-0.56\times1.5)^3=1-0.16^3=0.995904.
\label{eq:dh}
\end{equation}

The revised scalar is therefore
\begin{equation}
S_H=\frac{U}{D_H}=\frac{U}{0.995904}\in[0,1].
\label{eq:newscore}
\end{equation}

No cap is needed before division because $D_H$ is the maximum attainable value in
the defined input domain. An implementation may defensively evaluate
$\min(U,D_H)/D_H$.

Table~\ref{tab:cvss-pain-boundary} makes the boundary concrete. It compares MISS,
not a complete CVSS Environmental Score; a complete score would require the rest of
the CVSS vector and formula.

\begin{table}[H]
\centering
\small
\renewcommand{\arraystretch}{1.15}
\begin{tabularx}{\textwidth}{@{}P{0.30\textwidth}rrrY@{}}
\toprule
\textbf{Weighted-impact state} & \textbf{$U$} & \textbf{CVSS MISS} & \textbf{PAIN $S_H$} & \textbf{Effect} \\
\midrule
One H impact on one H requirement & 0.840000 & 0.840000 & 0.843455 & Below the cap; no information is clipped. \\
One H/H plus one H/M contribution & 0.929600 & 0.915000 & 0.933423 & CVSS clips $0.014600$; PAIN retains the compound increment. \\
Three H impacts on three H requirements & 0.995904 & 0.915000 & 1.000000 & CVSS clips $0.080904$; PAIN reaches its defined maximum. \\
\bottomrule
\end{tabularx}
\caption{Where the CVSS MISS path and PAIN normalization path diverge.}
\label{tab:cvss-pain-boundary}
\end{table}

\begin{cautionbox}
CVSS and PAIN use the same impact inputs, but they use them for different purposes.
CVSS stops the aggregate at $0.915$ and continues calculating the official CVSS
score. PAIN keeps the additional impact above $0.915$ and divides by $0.995904$ so
that strongly weighted dimensions remain distinguishable. This changes only the
PAIN classification. It does not recalculate or replace the published CVSS score or
vector. $U$ and $S_H$ must therefore be labeled as PAIN values, not CVSS scores.
\end{cautionbox}

\section{Deriving the thresholds}
\label{sec:thresholds}

\subsection{Deriving the Minimal/Narrow and Narrow/Disruptive boundaries}

The original $0.25$ and $0.55$ values were governed PAIN policy, not CVSS constants.
Merely multiplying them by $0.915/D_H$ would preserve their old raw-$U$ locations,
but it would not independently explain why those locations separate the PAIN words.
The revision instead applies the same scenario-first method used for Debilitating.

There is an important limit to this derivation. The C/I/A and CR/IR/AR lattice does
not encode the number of affected users, outage duration, or amount and type of
affected data. It therefore cannot reproduce every part of the FedRAMP word
definitions by arithmetic alone. The boundaries below are governed, standards-informed
proxies: CVSS supplies the discrete impact and requirement values; FIPS 199 supplies
the distinction between limited, serious, and severe or catastrophic adverse effect;
PAIN assigns those states to its operational words.

For a single dimension, the shorthand H/L means High technical impact on a Low
Security Requirement; L/H and H/M are read the same way. When terms are joined by
a plus sign, they occur on different C/I/A dimensions.

\paragraph{Narrow begins with a complete loss on a Low-requirement dimension.}
A High technical impact on a Low Security Requirement contributes
$0.56\times0.5=0.28$. Although the asset owner has determined that loss of the
property has limited adverse effect, the technical loss itself is complete. The
revision treats that state as Narrow rather than Minimal:
\begin{equation}
\boxed{\theta_1=\frac{0.28}{0.995904}\approx0.2811515969}.
\label{eq:theta1}
\end{equation}
An isolated Low impact on a Moderate requirement remains below the line:
$0.22/0.995904=0.2209048262$. Exhaustive enumeration confirms that no attainable
state lies between $0.2209048262$ and $0.2811515969$. With an inclusive comparison,
the complete-loss anchor itself is Narrow.

\paragraph{Disruptive begins with a High impact on a Moderate requirement.}
A Moderate Security Requirement represents a property whose loss can have a serious
adverse effect. A High technical impact directly aligned to that property therefore
provides the first unambiguous Disruptive anchor:
\begin{equation}
\boxed{\theta_2=\frac{0.56}{0.995904}\approx0.5623031939}.
\label{eq:theta2}
\end{equation}
The two closest disputed patterns remain Narrow. One H/L contribution combined
with one L/H contribution yields $S_H=0.5197288092$; two L/H contributions yield
$S_H=0.5533665896$. Neither contains a High technical impact on a Moderate-or-High
requirement. The latter is the greatest attainable state below the H/M anchor, and
no state lies between $0.5533665896$ and $0.5623031939$.

The boundary does not forbid lower-level contributions from becoming Disruptive.
Because the complement product remains continuous over all three dimensions, any
compound state whose aggregate reaches the H/M-equivalent value also crosses the
line. The rule is an arithmetic threshold, not a categorical H/M gate.

\subsection{Deriving Debilitating from the discrete state space}

For derivation only, define a \emph{compound catastrophic} state as one containing:

\begin{enumerate}[leftmargin=1.7em]
\item at least one High technical impact aligned with a High requirement; and
\item at least one additional High technical impact aligned with a Moderate-or-High
  requirement.
\end{enumerate}

This definition is not implemented as a gate. It partitions the finite input space
so that a scalar boundary can be derived.

There are $3^3=27$ possible C/I/A impact patterns and $3^3=27$ possible
CR/IR/AR patterns, for 729 combinations. Exhaustive enumeration establishes two
extrema.

The highest non-compound combination is one High impact on a High requirement plus
two Low impacts on High requirements:
\begin{align}
U_{\max,\mathrm{noncompound}}
  &=1-(1-0.84)(1-0.33)^2=0.928176,\\
S_{\max,\mathrm{noncompound}}
  &=\frac{0.928176}{0.995904}=0.9319934452.
\end{align}

The lowest compound catastrophic combination is one High impact on a High
requirement plus one High impact on a Moderate requirement, with no third impact:
\begin{align}
U_{\min,\mathrm{compound}}
  &=1-(1-0.84)(1-0.56)=0.929600,\\
S_{\min,\mathrm{compound}}
  &=\frac{0.929600}{0.995904}=0.9334233018.
\end{align}

The open interval between those values is non-empty. The selected threshold is
\begin{equation}
\boxed{\theta_3=0.933},
\label{eq:theta3}
\end{equation}
which satisfies
\[
0.9319934452 < 0.933 < 0.9334233018.
\]

Across all 729 discrete combinations, $S_H\ge0.933$ if and only if the combination
meets the compound catastrophic definition above. Seventy combinations qualify;
none of the 659 non-compound combinations cross the threshold, and no compound
combination falls below it.

\begin{decisionbox}
The final PAIN scalar and thresholds are:
\[
U=1-(1-C\mathrm{CR})(1-I\mathrm{IR})(1-A\mathrm{AR}),
\qquad
S_H=\frac{U}{0.995904},
\]
\[
(\theta_1,\theta_2,\theta_3)=
\left(\frac{0.28}{0.995904},\ \frac{0.56}{0.995904},\ 0.933\right)
\approx(0.2811515969,\ 0.5623031939,\ 0.933).
\]
Minimal is $S_H<\theta_1$; Narrow is $\theta_1\le S_H<\theta_2$;
Disruptive is $\theta_2\le S_H<\theta_3$; and Debilitating is
$S_H\ge\theta_3$. Implementations classify with full-precision $S_H$ and
thresholds. Rounded values are for display only. The parameter set must be versioned
as governed PAIN policy.
\end{decisionbox}

\section{Ceiling-constrained reachability}
\label{sec:ceiling-reachability}

The 729-state proof above describes the unconstrained Cartesian product of every
CVSS impact vector and every effective Security Requirements vector. A particular
deployment may expose only a subset of those states after its intended-use ceiling
is applied.

Let the uncapped archetype requirement be $R_x(o)$ and the intended-use ceiling be
$E_x(o)$. The PAIN calculation receives:
\[
R_x^{\mathrm{eff}}(o)=\min\big(R_x(o),E_x(o)\big).
\]
The ceiling derivation is upstream of this calibration; see \ceilingpaper{}. This
section examines only what the adopted PAIN arithmetic does with that input.

\subsection{Low- and Moderate-impact ceilings}

If a ceiling contains no High objective, then every effective requirement is at
most Moderate, regardless of the archetype. The largest possible pre-cap aggregate
occurs with High technical impact on all three Moderate requirements:
\begin{align}
U_{\max,\le M}
  &=1-(1-0.56)^3=0.914816,\\
S_{\max,\le M}
  &=\frac{0.914816}{0.995904}=0.918578\ldots
\end{align}
This is below the adopted Debilitating threshold:
\[
0.918578\ldots < 0.933.
\]
Therefore no state permitted by an all-Low/Moderate ceiling reaches Debilitating
under this calibration. All 70 Debilitating states in the unconstrained lattice
require at least one effective High requirement and are removed by such a ceiling.

\begin{cautionbox}
This is a consequence of the PAIN method's retained CVSS weights,
High-centered denominator, and $0.933$ threshold. NIST SP~800-60 and FIPS~199 do
not independently define this reachability result. A different PAIN aggregation
or threshold could classify the same effective requirements differently.
\end{cautionbox}

FIPS~199 separately establishes that any High system objective makes the
information system High overall. Combining that categorization rule with this
methodology-specific result means that a genuinely Low- or Moderate-impact
information system cannot reach the Debilitating word \emph{under the adopted PAIN
calibration}. The dimensional vector must still be retained: an overall High
system may be H/M/L, and the label High does not turn it into H/H/H.

\subsection{What a High ceiling does not prove}

A High objective makes Debilitating reachable in the deployment's state space; it
does not make every vulnerability Debilitating. The finding still needs:

\begin{enumerate}[leftmargin=1.7em]
\item a High technical impact aligned with an effective High requirement; and
\item another High technical impact aligned with an effective Moderate-or-High
  requirement.
\end{enumerate}

The archetype can also keep a dimension below the ceiling because the ceiling is
only a cap. Finally, multi-agency scope maps an already derived customer-effect word
to the N-level; it does not create the Debilitating word.

\section{Worked scenarios}
\label{sec:scenarios}

Tables~\ref{tab:lower-boundaries} and~\ref{tab:scenarios} apply the final arithmetic
to boundary and ordinary cases. Requirement vectors are explicit effective inputs:
the governed asset catalog supplies the uncapped archetype and any intended-use
ceiling caps it before these calculations. Results assume a single-agency resource;
the separate agency-scope mapping can promote qualifying words to N4 or N5.

\begin{table}[H]
\centering
\footnotesize
\renewcommand{\arraystretch}{1.16}
\setlength{\tabcolsep}{4pt}
\begin{tabularx}{0.98\textwidth}{@{}P{0.21\textwidth}P{0.08\textwidth}P{0.12\textwidth}rrY@{}}
\toprule
\textbf{Boundary scenario} & \textbf{C/I/A} & \textbf{CR/IR/AR} & \textbf{$U$} & \textbf{$S_H$} & \textbf{Result and reason} \\
\midrule
Isolated limited effect & L/N/N & M/M/M & 0.220000 & 0.220905 & N1. Low impact on a Moderate requirement remains Minimal. \\
Complete loss, Low requirement & H/N/N & L/M/M & 0.280000 & 0.281152 & N2. First Narrow anchor. \\
Crossed mixed pair & H/L/N & L/H/M & 0.517600 & 0.519729 & N2. No direct serious alignment. \\
Two limited effects on High requirements & L/L/N & H/H/M & 0.551100 & 0.553367 & N2. Greatest attainable state below H/M. \\
Direct serious alignment & H/N/N & M/M/M & 0.560000 & 0.562303 & N3. First Disruptive anchor. \\
\bottomrule
\end{tabularx}
\caption{Strict lower-boundary scenarios. Threshold comparisons are inclusive.}
\label{tab:lower-boundaries}
\end{table}

\begin{table}[H]
\centering
\footnotesize
\renewcommand{\arraystretch}{1.18}
\setlength{\tabcolsep}{4pt}
\begin{tabularx}{0.98\textwidth}{@{}P{0.20\textwidth}P{0.08\textwidth}P{0.12\textwidth}rrY@{}}
\toprule
\textbf{Scenario} & \textbf{C/I/A} & \textbf{CR/IR/AR} & \textbf{$U$} & \textbf{$S_H$} & \textbf{Result and reason} \\
\midrule
Availability-only orchestrator & N/N/H & M/M/H & 0.840000 & 0.843455 & N3. Catastrophic availability is isolated. \\
Integrity and availability on orchestrator & N/H/H & M/M/H & 0.929600 & 0.933423 & N4. Minimum compound catastrophic case. \\
Complete orchestrator compromise & H/H/H & M/M/H & 0.969024 & 0.973010 & N4. Catastrophic availability plus serious C/I impact. \\
Confidentiality and integrity on data backbone & H/H/N & H/H/M & 0.974400 & 0.978408 & N4. Two catastrophic dimensions. \\
Complete app-tier compromise & H/H/H & M/M/M & 0.914816 & 0.918578 & N3. Broad serious impact, but no catastrophic requirement alignment. \\
Worst non-compound boundary & H/L/L & H/H/H & 0.928176 & 0.931993 & N3. Low-impact accumulation remains below the line. \\
Low availability on orchestrator & N/N/L & M/M/H & 0.330000 & 0.331358 & N2. Partial availability impact remains Narrow. \\
\bottomrule
\end{tabularx}
\caption{Ordinary and Debilitating-boundary behavior under High-centered normalization.}
\label{tab:scenarios}
\end{table}

The first two rows of Table~\ref{tab:scenarios} are the central upper-tier
separation. A single total loss on a
catastrophic dimension is Disruptive. Adding a High impact on a second serious
dimension crosses the derived line. The fifth row is an explicit policy decision:
complete CIA compromise of an all-Moderate asset remains N3 because the governed
requirements describe every loss as serious rather than catastrophic.

\paragraph{Concrete finding example.}
Dataset A contains a denial-of-service finding with $C{:}N/I{:}N/A{:}H$ on a
resource governed by an \code{MMH} requirement vector. $U=0.84$ and
$S_H=0.843455$, producing N3. The asset can retain the correct High Availability
Requirement without creating an N4 deadline from a one-dimensional impact.

\section{Empirical backtest}
\label{sec:backtest}

The candidate methods were replayed against two anonymized evaluation reports
generated on July 21, 2026. Dataset A contained 2,120 findings and 10,231
affected-resource observations. Dataset B contained 3,691 findings and
60,817 affected-resource observations. The same CVE may affect multiple resources;
the two units are therefore reported separately.

All candidate runs used the same adopted archetype catalog and the existing
engine's impact-source precedence, including CVSS v3/v4 vector parsing, qualitative
severity fallback, and technical-impact lifting where recorded. The stored-current
rows retain the reports' original policy and provide operational context; they are
not an isolated single-variable control.

\begin{table}[H]
\centering
\small
\renewcommand{\arraystretch}{1.15}
\begin{tabular}{@{}lrr@{}}
\toprule
\textbf{Method} & \textbf{Dataset A N4} & \textbf{Dataset B N4} \\
\midrule
Stored scan output & 40.56\% & 63.95\% \\
Adopted catalog baseline; original $D=0.915$, $\theta_3=0.80$ & 26.62\% & 62.22\% \\
Raise $\theta_3$ to $0.92$ only & 15.66\% & 29.54\% \\
Change High multiplier to $1.25$ only & 17.49\% & 29.76\% \\
Strict two-H/H gate & 6.61\% & 20.39\% \\
Compound gate & 11.60\% & 21.22\% \\
\textbf{Final: $D_H=0.995904$, $\theta_3=0.933$} & \textbf{11.60\%} & \textbf{21.22\%} \\
\bottomrule
\end{tabular}
\caption{N4 share of affected-resource observations.}
\label{tab:affected-n4}
\end{table}

The final arithmetic selected exactly the same N4 affected-resource observations
as the compound gate in both reports. With the lower thresholds independently
anchored to the H/L and H/M scenarios, the complete distributions are shown in
Table~\ref{tab:distribution}.

\begin{table}[H]
\centering
\small
\renewcommand{\arraystretch}{1.15}
\begin{tabular}{@{}lrrrr@{}}
\toprule
\textbf{Scan and policy} & \textbf{N1} & \textbf{N2} & \textbf{N3} & \textbf{N4} \\
\midrule
Dataset A stored current & 11.37\% & 28.47\% & 19.60\% & 40.56\% \\
Dataset A final & 16.12\% & 27.08\% & 45.20\% & \textbf{11.60\%} \\
Dataset B stored current & 1.50\% & 28.29\% & 6.25\% & 63.95\% \\
Dataset B final & 2.21\% & 30.41\% & 46.16\% & \textbf{21.22\%} \\
\bottomrule
\end{tabular}
\caption{Affected-resource PAIN distributions under the final calibration.}
\label{tab:distribution}
\end{table}

At the finding level, taking the worst affected resource per finding, N4 moved
from 55.80\% to 18.16\% in Dataset A and from 50.50\% to 17.83\% in Dataset B.

\subsection{The threshold remains reachable}

The purpose of $0.933$ is not to make N4 exceptional or unattainable. It is crossed
by the common two-High impact shapes $H/H/N$, $H/N/H$, and $N/H/H$, as well as
$H/H/H$, whenever the requirements provide the necessary catastrophic and serious
alignment.

\begin{plain}
$0.933$ does not mean that a vulnerability must cause ``93.3 percent damage,'' nor
does it represent the 93.3rd percentile of observed findings. The PAIN scalar is an
ordered index over discrete CVSS-derived states. The first qualifying compound state
lands at $0.933423$, immediately above the line, while ordinary two-High and
three-High scenarios land from that point through $1.0$. A vulnerability does not
need three High impacts, three High requirements, multi-agency scope, or a specially
constructed vector to reach Debilitating.
\end{plain}

Dataset A retained 1,187 N4 affected-resource observations across multiple archetypes:
701 \emph{data-backbone}, 177 \emph{orchestrator}, 125
\emph{config-actuation}, 102 \emph{identity-secrets}, 49
\emph{security-tooling}, and 33 across other roles. Its most frequent qualifying
patterns included $HHH/HHM$, $HHH/MMH$, $HHN/HHM$, $HHH/MHM$, and $HHH/HMM$.

Dataset B retained 12,903 N4 observations, but 11,622 of them were
\code{unclassified} and therefore used the conservative $HHH$ fail-safe. This is an
important interpretation constraint: the remaining Dataset B N4 concentration is
primarily a classification-coverage issue, not evidence that $0.933$ is too low.

At the worst-resource-per-finding level, 385 Dataset A findings (18.16\%) and 658
Dataset B findings (17.83\%) remained N4. A top tier reached by approximately one in
six findings in both snapshots is plainly operational rather than theoretical. The
method nevertheless removes the one-dimensional cases that made N4 cease to signal
the worst compound consequences.

\subsection{What the backtest does and does not prove}

The backtest demonstrates that the threshold is operationally reachable, does not
collapse N4 to a handful of exotic vectors, and produces the same top-tier set as
the independently stated compound predicate. It does not prove that 11.60\% or
21.22\% is an ideal universal N4 rate. The reports are two deployment snapshots,
not a representative sample of all federal cloud systems. The mathematical state
boundary supplies the policy justification; the scans test its consequences.

\section{Why the final method is preferable}

\begin{table}[H]
\centering
\footnotesize
\renewcommand{\arraystretch}{1.18}
\setlength{\tabcolsep}{5pt}
\begin{tabularx}{\textwidth}{@{}P{0.19\textwidth}YYYY@{}}
\toprule
\textbf{Approach} & \textbf{CVSS constants} & \textbf{Tagging burden} & \textbf{Dimensionality} & \textbf{Boundary rationale} \\
\midrule
Asset sub-profiles & Preserved & High & Preserved & Hidden in classification choices \\
$1.25$ multiplier & Changed & Low & Preserved & Indirect \\
Threshold-only & Preserved & Low & Preserved & Weak unless separately derived \\
Strict gate & Preserved & Low & Reduced to H/H count & Too restrictive \\
Compound gate & Preserved & Low & Reduced at the final branch & Explicit but adds a second classifier \\
High-centered arithmetic & Preserved & Low & Preserved continuously & Derived from the full discrete state space \\
\bottomrule
\end{tabularx}
\caption{Decision summary. ``CVSS constants preserved'' does not mean the final
PAIN scalar is a conformant CVSS score; the normalization is explicitly PAIN-specific.}
\end{table}

The final method is preferable because it obtains the compound-gate result without
executing a gate. It keeps Low impacts and all three requirement levels inside one
continuous calculation, retains the official metric weights, and makes the top
threshold reproducible from the input lattice. It also keeps the operator-facing
model simple: one role-based archetype supplies the reusable architectural vector,
and one scope-resolved intended-use ceiling supplies the federal consequence cap.
Neither changes per finding.

\section{Implementation and governance}

\subsection{Normative parameter set}

An implementation of this calibration should govern the following values together:

\begin{verbatim}
impactWeights:
  high: 0.56
  low: 0.22
  none: 0.0

requirementWeights:
  high: 1.5
  medium: 1.0
  low: 0.5

normalizationDenominator: 0.995904

wordThresholds:
  narrow: 0.28115159694107
  disruptive: 0.56230319388214
  debilitating: 0.933
\end{verbatim}

The denominator and thresholds are one calibration set. The first two thresholds
should be constructed from the exact expressions $0.28/D_H$ and $0.56/D_H$ when the
implementation permits it. The decimal serializations above are deliberately
truncated rather than rounded upward so that the H/L and H/M anchors remain inside
their intended bands under an inclusive comparison. Changing a metric weight, adding
an impact level, or changing the aggregation invalidates the boundary proof and
requires re-enumeration.

\subsection{Precision and reproducibility}

The separation between the highest non-compound and lowest compound states is
approximately $0.00143$. Implementations therefore \rfckw{must}:

\begin{itemize}[leftmargin=1.5em]
\item classify using unrounded internal values;
\item construct $\theta_1$ and $\theta_2$ from $0.28/D_H$ and $0.56/D_H$, or use
  decimal serializations that do not round above those exact anchors;
\item compare against the configured full-precision threshold;
\item display enough decimal places to explain a boundary decision;
\item record $U$, $D_H$, $S_H$, the thresholds, resolved C/I/A, CR/IR/AR, and the
  policy version in audit output; and
\item include boundary unit tests for the predecessor/anchor pairs
  $0.2209048262/0.2811515969$, $0.5533665896/0.5623031939$, and
  $0.9319934452/0.9334233018$.
\end{itemize}

\subsection{Migration}

The change should be released as a new PAIN methodology version. Historical
findings should retain the policy version under which their deadline was assigned.
Active findings should be recomputed under an explicitly approved transition rule;
the organization must acknowledge that N4-to-N3 changes extend remediation
deadlines. The scoring calibration should be versioned and reviewed independently
from any asset-catalog revision.

The default for an unclassified asset remains conservative $HHH$. The final method
reduces single-dimension findings on that default to N3, but two- and three-High
impact findings readily remain N4. Classification coverage should be monitored so
that the fail-safe does not become the dominant production archetype.

\section{Limitations and future work}

\begin{itemize}[leftmargin=1.5em]
\item The exact $0.933$ separation depends on the discrete
  $\{N,L,H\}\times\{L,M,H\}$ domain and the retained CVSS v3.1 weights. It must be
  re-derived if those inputs change.
\item The finding that an all-Low/Moderate ceiling cannot reach Debilitating is a
  consequence of this calibration, not an independent NIST or FedRAMP rule. It
  must be re-evaluated if the aggregation, weights, denominator, or thresholds
  change.
\item Native CVSS v4 subsequent-system impacts and Safety are not included in the
  three-dimension proof. A future PAIN version may add them, but doing so requires a
  new denominator and threshold analysis rather than silently extending this one.
\item The two backtests establish observed consequences, not statistical
  generalizability. Additional production datasets and analyst-rated scenario sets
  should be evaluated before broad standardization.
\item A scalar cannot replace architecture review. Archetype assignment, Modified
  impact evidence, intended-use ceiling, agency scope, and vulnerability
  applicability remain governed inputs.
\item The threshold's narrow mathematical gap is safe for deterministic software
  using full precision but unsuitable for manual classification from rounded
  dashboard values.
\end{itemize}

\section{Conclusion}

The original PAIN method correctly joined vulnerability impact with asset-specific
security requirements, but its Medium-centered denominator saturated too early for
High-requirement assets. The resulting single-dimension cliff made Debilitating too
common and created pressure to solve a scoring problem through asset
underclassification.

The review considered richer tagging, a lower High multiplier, threshold-only
tuning, and categorical gates. Each could reduce N4, but each either increased
operator complexity, abandoned a CVSS-defined value, obscured the semantic reason
for the threshold, or discarded dimensional information.

High-centered normalization resolves the underlying mismatch. The new denominator
is derived directly from the retained CVSS weights. Narrow begins at a complete
technical loss on a Low-requirement dimension; Disruptive begins at a High impact
on a Moderate-requirement dimension. The Debilitating threshold lies between two
provable extrema in the complete discrete state space and remains readily reachable
by ordinary two-High impact scenarios. The empirical scans then validate the
operational consequences rather than dictate the boundaries.

The resulting rule is concise:
\[
\boxed{
S_H=\frac{1-(1-C\mathrm{CR})(1-I\mathrm{IR})(1-A\mathrm{AR})}{0.995904}
}
\]
with
\[
\boxed{
(\theta_1,\theta_2,\theta_3)
=\left(\frac{0.28}{0.995904},\ \frac{0.56}{0.995904},\ 0.933\right).
}
\]

This preserves the CVSS-derived evidence chain while giving PAIN a more defensible
meaning at its most consequential boundary: N4 and N5 are reserved for impact that
is not only severe, but compound and catastrophic in the affected environment.
When an intended-use ceiling has no High objective, the adopted formula makes that
boundary unreachable; this is a transparent property of the PAIN calibration,
not a categorization rule imported from NIST or FedRAMP.

\clearpage
\section*{References}
\addcontentsline{toc}{section}{References}

{\small
\begin{itemize}[leftmargin=3.2em,itemsep=0.3em,topsep=0.4em]
\item[{[PAIN]}] M.~Venne, \emph{A Deterministic, CVSS--Environmental Method for
  FedRAMP VDR/VER Vulnerability Prioritization}, stackArmor, July 2026.
  \url{https://stackarmor.github.io/rfc-fedramp-vdr/vdr-pain-cvss.pdf}.
\item[{[CEILING]}] M.~Venne, \emph{Before the PAIN Equation: Deriving
  Security-Requirements Ceilings from Intended Federal Information Types},
  stackArmor, July 2026.
  \url{https://stackarmor.github.io/rfc-fedramp-vdr/vdr-pain-security-requirements.pdf}.
\item[{[CVSS31]}] FIRST, \emph{Common Vulnerability Scoring System v3.1:
  Specification Document}. Environmental Security Requirements, MISS formula, and
  metric values. \url{https://www.first.org/cvss/v3.1/specification-document}.
\item[{[CVSS40]}] FIRST, \emph{Common Vulnerability Scoring System v4.0:
  Specification Document}. Environmental metrics and EQ3/EQ6 MacroVectors.
  \url{https://www.first.org/cvss/v4.0/specification-document}.
\item[{[CVSS40-UG]}] FIRST, \emph{CVSS v4.0 User Guide}. Expert comparison,
  equivalence-set, and interpolation methodology.
  \url{https://www.first.org/cvss/user-guide.html}.
\item[{[FIPS199]}] National Institute of Standards and Technology,
  \emph{Standards for Security Categorization of Federal Information and
  Information Systems}, FIPS PUB 199.
  \url{https://csrc.nist.gov/pubs/fips/199/final}.
\item[{[SP800-60]}] National Institute of Standards and Technology,
  \emph{Guide for Mapping Types of Information and Information Systems to Security
  Categories}, SP 800-60 Volumes I and II, Revision 1.
  \url{https://csrc.nist.gov/pubs/sp/800/60/r1/final}.
\item[{[FEDRAMP]}] FedRAMP, \emph{Potential Agency Impact Definitions},
  Notice 0012, and \emph{Vulnerability Evaluation and Reporting}, Consolidated
  Rules 2026.
  \url{https://www.fedramp.gov/notices/0012/};
  \url{https://www.fedramp.gov/2026/reference/rev5/b/vulnerability-evaluation-and-reporting/}.
\end{itemize}
}

\subsection*{Reproducibility record}

The discrete-state verification enumerates the Cartesian product of three impact
values for each of C/I/A and three requirement values for each of CR/IR/AR. A
conforming test run over the resulting 729 states must return:

\begin{center}
\small
\begin{tabular}{@{}lr@{}}
\toprule
\textbf{Check} & \textbf{Expected result} \\
\midrule
$D_H$ & 0.995904000000 \\
Predecessor to Narrow & 0.22090482616798 \\
Narrow anchor (H/L) & 0.28115159694107 \\
Predecessor to Disruptive & 0.55336658955080 \\
Disruptive anchor (H/M) & 0.56230319388214 \\
Maximum non-compound $S_H$ & 0.93199344515134 \\
Minimum compound $S_H$ & 0.93342330184435 \\
States at or above $0.933$ & 70 \\
State counts Minimal/Narrow/Disruptive/Debilitating & 90 / 218 / 351 / 70 \\
Predicate/threshold mismatches & 0 \\
Maximum $S_H$ with every requirement $\le M$ & 0.918578497526 \\
Debilitating states after an all-Low/Moderate ceiling & 0 \\
\bottomrule
\end{tabular}
\end{center}

The empirical run used the two July 21, 2026 evaluation exports described in
Section~\ref{sec:backtest}: 2,120 findings and 10,231 affected-resource observations
for Dataset A; 3,691 findings and 60,817 observations for Dataset B. Dataset names,
resource names, and source identifiers are intentionally omitted from this public
paper. Authorized reviewers may reproduce the run using the controlled internal
evidence package.

\end{document}
